\documentclass[11pt]{article}
\usepackage[margin=1in]{geometry}
\usepackage{booktabs}
\usepackage{hyperref}
\usepackage{siunitx}
\usepackage{xcolor}
\usepackage{parskip}

\title{LUCID Slot~040 --- Replication Report\\
\large Park \emph{et al.} (2022) --- New Damage Model for Direct Damage in
Biomolecular Systems (Geant4-DNA, pBR322)}
\author{OpenClaw Replication Subagent (Argo Opus 4.7, free endpoint)}
\date{2026-06-22 (report), 2026-07-05 (backfill \& critique)}

\begin{document}
\maketitle

\section{Paper}
Park~J., Jung~K.-W., Kim~M.K., Gwon~H.-J., Jung~J.-H.\ (2022).
``New damage model for simulating radiation-induced direct damage to
biomolecular systems and experimental validation using pBR322 plasmid.''
\emph{Scientific Reports} \textbf{12}, 11345.
\href{https://doi.org/10.1038/s41598-022-15521-y}{doi:10.1038/s41598-022-15521-y}.
Korea Atomic Energy Research Institute (KAERI); Geant4-DNA v10.7p01.

\section{Verdict (four-tier, cross-checked)}
\textbf{Queue verdict:} REPLICATED. \textbf{This report's substantive
verdict:} \emph{Partial --- Analytical Complete}. The two verdicts are
consistent \emph{if} ``replicated'' is read as ``every analytically
tractable claim reproduced''; the MC-dependent headline claim
(\SI{14.2}{\percent} mean error against gel-electrophoresis SSB/DSB data,
seven LET points, Geant4-DNA v10.7p01 on 5{,}400 plasmids in a
\SI{3}{\micro\meter} sphere) was \textbf{not} reproduced here and is not
independently reproducible from published information alone. We preserve
the queue verdict as \textbf{REPLICATED} but flag that the label is
misleadingly optimistic --- see \S\ref{sec:critique}.

\begin{center}
\begin{tabular}{lll}
\toprule
Tier & Definition & This work \\
\midrule
Full        & All quantitative claims independently reproduced & --- \\
\textbf{Partial}    & \textbf{Analytical scope fully closed; MC scope out of reach} & \(\checkmark\) \\
Conceptual  & Method described; numbers diverge & --- \\
Not reproduced & Numbers cannot be regenerated from paper & --- \\
\bottomrule
\end{tabular}
\end{center}

\textbf{Coverage 6/10 \quad Agreement 9/10} (within analytical scope).

\section{What was actually reproduced (analytical)}
\begin{itemize}
  \item Coarse-grained (CG) bead radii from group volumes: phosphate
        \SI{2.285}{\angstrom} (paper \SI{2.3}{\angstrom}), deoxyribose
        \SI{2.717}{\angstrom} (\SI{2.7}{\angstrom}), base
        \SI{2.917}{\angstrom} (\SI{2.9}{\angstrom}). Agreement
        \(<\SI{0.02}{\angstrom}\) throughout.
  \item Table 2 phosphate PO$_3$ total binding energy
        \(-12.3566\) eV vs paper \(-12.3562\) eV (\SI{0.4}{\milli\electronvolt}).
  \item McMahon--Currell mass-conservation residual $\sim 6\times10^{-17}$
        (machine precision).
  \item Back-fit of $\mu$ from OC$(D)$ curve for both $^{60}$Co and
        \SI{1}{MeV}~$e^-$ recovers input to $\leq 0.4\%$.
  \item Eq.~(8) mean-percentage-error formula itself: exact.
\end{itemize}

\section{What was NOT reproduced}
Figures~3--5 (SSB/DSB yields vs LET, plasmid vs linear DNA, comparison
to TOPAS-nBio) and the headline \SI{14.2}{\percent} mean-error claim all
depend on a Geant4-DNA v10.7p01 MC run over seven beams and
5{,}400 plasmids, requiring the authors' unreleased application code and
GPU-scale compute. None of that is analytically tractable.

\section{Bug discovery: paper Eq.~(3) is misprinted}
The typeset Morse expression
\(U_1(r) = D_e\{1 - e^{-2\alpha(r-r_e)} - 2e^{-\alpha(r-r_e)}\}\)
yields $U_1(r_e) = -2 D_e$ and inflates the phosphate CG total to
$\sim -34.0$~eV --- $2.75\times$ the paper's own Table~2 sum. The
\emph{standard} Morse form $D_e(1-e^{-\alpha(r-r_e)})^2 - D_e$ (which
gives $U_1(r_e)=-D_e$) with the same Table~1 parameters reproduces
every Table~2 row to $<\SI{7}{\milli\electronvolt}$. Table~2 was
therefore produced with the standard form; only the typeset Eq.~(3) is
wrong. Useful errata.

\section{Genuine critique}\label{sec:critique}
The following weaknesses matter for anyone deciding whether to trust or
build on this model.

\paragraph{1. The direct/indirect split is not independently derived.}
The paper's central sales-pitch is a ``new damage model'' that replaces
empirically tuned threshold energies with CG-potential threshold
energies for the \emph{direct} channel. But the paper does not perform
an independent chemistry-side calibration for the \emph{indirect}
(radical-mediated) channel; the McMahon $\phi$ parameter that soaks up
indirect damage is fit to totals, not derived from OH-radical yields or
diffusion kinetics. In our back-fits the OC$(D)$ curve recovered $\mu$
to $<0.4\%$ but $\phi$ only to $\pm 46\%$ ($^{60}$Co) and $\pm 42\%$
(\SI{1}{MeV}~$e^-$). That is a known weak-constraint problem
of OC alone --- and the paper does not acknowledge it. A more honest
framing would say: the direct channel is now first-principles, but the
overall model is still one adjustable parameter ($\phi$) short of
identifiability from OC data alone.

\paragraph{2. No scavenger-modulation or OER validation.}
The obvious way to prove a \emph{direct} damage model is to vary the
radical scavenger concentration (DMSO, glycerol, Tris) or the oxygen
tension and show that the direct component --- and only the direct
component --- is invariant. The paper does neither. All validation data
are unmodulated $^{60}$Co and \SI{1}{MeV}~$e^-$ irradiation of dried
pBR322 in air. Dried plasmid samples suppress indirect damage, but do
not eliminate the bound-water shell. Without a scavenger sweep or an
OER measurement across the seven LET points, the model's central
claim --- that it faithfully partitions direct vs indirect --- is not
tested; only totals are tested.

\paragraph{3. No cross-code comparison beyond a single TOPAS-nBio panel.}
Fig.~5b compares to TOPAS-nBio, which shares Geant4-DNA physics.
No comparison is made to TRAX, KURBUC, or PARTRAC, all of which have
independent track-structure implementations and would provide a genuine
cross-check on the CG-potential threshold rather than an intra-family
consistency check.

\paragraph{4. Deoxyribose 30.5~eV cannot be reconstructed from the paper.}
Table~2 itemises only the phosphate PO$_3$ bead. The deoxyribose
30.5~eV threshold appears in the main text as a final number and is
never broken down in the article. Our closed-form estimate using the
canonical furanose skeleton (3 C--C + 3 C--O + 1 O$\cdots$O non-bonded)
yields 22.12~eV --- 8~eV short. The missing 8~eV very plausibly comes
from additional non-bonded pairs itemised in Supplementary Fig.~S2, but
S2 is behind a Nature supplement URL we did not retrieve. Anyone
reimplementing from the article alone will land at $\sim 22$~eV, not
$\sim 30$~eV, and the discrepancy propagates directly into
predicted deoxyribose-lesion yields.

\paragraph{5. Data-availability statement is functionally closed.}
``Available from the corresponding author through a reasonable written
request'' means no public code, no CG geometry file in ``DNAFabric
format,'' no SSB/DSB clustering algorithm, and no gel-band-intensity
CSV. All five reproducibility blockers documented in \S6 of the source
REPORT.md are consequences of this closed statement. In LUCID-corpus
terms this is a mid-tier reproducibility posture: better than
``code available on request'' meaning nothing, worse than any real
open-code release. Every subsequent reproduction attempt --- ours
included --- will hit the same wall.

\section{Files delivered (this backfill)}
\begin{itemize}
  \item \texttt{report/REPORT.md} --- primary analytical narrative (unchanged).
  \item \texttt{report/REPORT.tex} --- this file.
  \item \texttt{report/open\_questions.json} + \texttt{\_section.tex} --- 5 truly-open questions.
  \item \texttt{report/workflow.md} --- what was actually run.
  \item \texttt{report/artifacts\_summary.md} --- pointer index of every file.
  \item \texttt{report/failure\_analysis.md} --- honest critique of both paper \emph{and} this replication.
  \item \texttt{extraction/nougat.mmd} --- stub (source PDF was extracted via pdftotext into \texttt{ocr/raw\_layout.txt}; no Nougat run this session).
\end{itemize}

\end{document}
